\section{Regulatory and Governance Interpretation}

FL-BSA provides technical evidence for human review. It does not decide whether a lender has met a
legal obligation. Applicable law depends on jurisdiction, product, use, data provenance, protected
classes, decision process, and current regulatory interpretation. Reviewers should involve qualified
legal and compliance owners.

\subsection{Fair-lending review}

AIR and SRG can help surface selection-rate disparities under ECOA/Regulation B governance. The 0.80
screen is a predeclared internal heuristic; employment-selection guidance is historical context,
not a fair-lending safe harbor \citep{eeoc_uniform_guidelines}. A result outside the screen is a
review trigger. A result within it is not proof of nondiscrimination. Causal analysis, legitimate
business-need assessment, less-discriminatory alternatives, data quality, sample coverage, and
adverse-action processes remain outside this synthetic characterization. The current primary rule
bars use of prohibited bases in credit evaluation subject to its stated exceptions; this paper does
not translate a statistical screen into a conclusion under that rule \citep{ecfr_regb_1002_6}.
Effective 21 July 2026, the CFPB's final rule removed the Regulation B effects-test provision and
states that ECOA does not recognize disparate-impact liability \citep{cfpb_regb_2026}. This
engineering screen neither applies nor substitutes for a legal liability test; other laws,
jurisdictions, and authorities may differ.

\subsection{Model-risk governance}

The release-bound identities, limitations, independent branches, deterministic inputs, and
machine-readable evidence support model inventory, validation, challenge, change control, and
retention practices of the kind addressed by model-risk guidance \citep{sr26_2}. The artifact does
not replace independent validation. The advisory utility weakness and unavailable metrics are
review findings that should remain visible in any evaluation record.

\subsection{EU AI Act and FCA Consumer Duty}

For an EU AI Act assessment, the evidence can support technical documentation, logging,
traceability, data-governance discussion, accuracy/robustness review, and human-oversight design
\citep{eu_ai_act}. Classification, conformity assessment, deployer/provider responsibilities, and
legal conclusions require a separate facts-and-law analysis.

For FCA Consumer Duty governance, disparity screens and evidence retention can contribute to
monitoring customer outcomes, vulnerable-customer impacts, foreseeable harm, and governance
challenge \citep{fca_consumer_duty_2026}. This synthetic fixture is not evidence of outcomes for a
firm's actual customers and does not establish compliance with the Duty.

\subsection{Required human decisions}

Before broader use, an accountable organization should decide at least: the lawful basis and data
scope; relevant protected groups and intersectional slices; reference-group policy; minimum sample
sizes; uncertainty and multiplicity treatment; action thresholds; escalation ownership; model and
data change triggers; retention; and the boundary between technical screening and legal judgment.
